
\documentclass{article}
\usepackage{colortbl}
\usepackage{makecell}
\usepackage{multirow}
\usepackage{supertabular}

\begin{document}

\newcounter{utterance}

\centering \large Interaction Transcript for game `referencegame', experiment `number\_grids', episode 5 with reflexion{-}llama{-}3.1{-}8b{-}instruct.
\vspace{24pt}

{ \footnotesize  \setcounter{utterance}{1}
\setlength{\tabcolsep}{0pt}
\begin{supertabular}{c@{$\;$}|p{.15\linewidth}@{}p{.15\linewidth}p{.15\linewidth}p{.15\linewidth}p{.15\linewidth}p{.15\linewidth}}
    \# & $\;$A & \multicolumn{4}{c}{Game Master} & $\;\:$B\\
    \hline

    \theutterance \stepcounter{utterance}  
    & & \multicolumn{4}{p{0.6\linewidth}}{
        \cellcolor[rgb]{0.9,0.9,0.9}{
            \makecell[{{p{\linewidth}}}]{
                \texttt{\tiny{[P1$\langle$GM]}}
                \texttt{You have attempted this same task before and did not succeed. These are the plans you wrote after those attempts.} \\
\texttt{Trial \#1: New plan: I failed to distinguish the target grid from the distractor grids because I focused on the rows of X's, which are present in all three grids. Next time, I will look for a unique feature that is present in the target grid but not in the distractor grids. I will examine the columns and the positions of the X's more closely to find a distinguishing characteristic. For example, I might look for a column that has a unique pattern of X's and □'s, or a specific position where the X's are located.} \\
\\ 
\\ 
\texttt{You are given three grids, where each of them is 5 by 5 in size.} \\
\texttt{Grids have empty cells marked with "▢" and filled cells marked with "X".} \\
\texttt{Your task is to generate a referring expression that best describes the target grid while distinguishing it from the two other distractor grids.} \\
\texttt{The first grid is the target grid, and the following two grids are the distractors.} \\
\\ 
\texttt{Target grid:} \\
\\ 
\texttt{□ X X X X} \\
\texttt{□ X □ □ X} \\
\texttt{□ X X X X} \\
\texttt{□ X □ □ X} \\
\texttt{□ X X X X} \\
\\ 
\texttt{Distractor grid 1:} \\
\\ 
\texttt{□ X X X X} \\
\texttt{□ X □ □ X} \\
\texttt{□ X X X X} \\
\texttt{□ □ □ □ □} \\
\texttt{□ X X X X} \\
\\ 
\texttt{Distractor grid 2:} \\
\\ 
\texttt{□ X X X X} \\
\texttt{□ X □ □ X} \\
\texttt{□ X □ □ X} \\
\texttt{□ X □ □ X} \\
\texttt{□ X X X X} \\
\\ 
\texttt{Instruction: Describe the target grid.} \\
\texttt{Generate the referring expression starting with the tag "Expression: " for the given target grid. Omit any other text.} \\
            }
        }
    }
    & \\ \\

    \theutterance \stepcounter{utterance}  
    & \multicolumn{4}{p{0.6\linewidth}}{
        \cellcolor[rgb]{0.9,0.9,0.9}{
            \makecell[{{p{\linewidth}}}]{
                \texttt{\tiny{[P1$\rangle$GM]}}
                \texttt{Expression: The grid that has a row of five X's in the bottom row.} \\
            }
        }
    }
    & & \\ \\

    \theutterance \stepcounter{utterance}  
    & & & \multicolumn{2}{p{0.3\linewidth}}{
        \cellcolor[rgb]{0.9,0.9,0.9}{
            \makecell[{{p{\linewidth}}}]{
                \texttt{\tiny{[GM$|$GM]}}
                \texttt{Expression: The grid that has a row of five X's in the bottom row.} \\
            }
        }
    }
    & & \\ \\

    \theutterance \stepcounter{utterance}  
    & & \multicolumn{4}{p{0.6\linewidth}}{
        \cellcolor[rgb]{0.9,0.9,0.9}{
            \makecell[{{p{\linewidth}}}]{
                \texttt{\tiny{[GM$\rangle$P2]}}
                \texttt{You have attempted this same task before and did not succeed. These are the plans you wrote after those attempts.} \\
\texttt{Trial \#1: New plan: I failed to match the given expression because I didn't carefully read it. The expression states that the grid should have a row of five X's in the top row and a row of five X's in the bottom row. However, I selected the first grid, which only has a row of five X's in the top row. Next time, I will carefully read the expression and look for the additional condition of a row of five X's in the bottom row. I will select the grid that meets both conditions.} \\
\\ 
\\ 
\texttt{You are given three grids, where each of them is 5 by 5 in size.} \\
\texttt{Grids have empty cells marked with "▢" and filled cells marked with "X".} \\
\texttt{You are also given a referring expression that describes one of the given grids.} \\
\texttt{Your task is to select the grid that matches the given referring expression.} \\
\texttt{Generate only the number (in text) of the grid that matches the given expression by selecting first, second, or third.} \\
\\ 
\texttt{First:} \\
\\ 
\texttt{□ X X X X} \\
\texttt{□ X □ □ X} \\
\texttt{□ X □ □ X} \\
\texttt{□ X □ □ X} \\
\texttt{□ X X X X} \\
\\ 
\texttt{Second:} \\
\\ 
\texttt{□ X X X X} \\
\texttt{□ X □ □ X} \\
\texttt{□ X X X X} \\
\texttt{□ □ □ □ □} \\
\texttt{□ X X X X} \\
\\ 
\texttt{Third:} \\
\\ 
\texttt{□ X X X X} \\
\texttt{□ X □ □ X} \\
\texttt{□ X X X X} \\
\texttt{□ X □ □ X} \\
\texttt{□ X X X X} \\
\\ 
\texttt{Expression: The grid that has a row of five X's in the bottom row.} \\
\texttt{Question: Which grid does the expression refer to?} \\
\texttt{Start with the tag "Answer: ", followed by your selection. Omit any other text.} \\
            }
        }
    }
    & \\ \\

    \theutterance \stepcounter{utterance}  
    & & & \multicolumn{4}{p{0.6\linewidth}}{
        \cellcolor[rgb]{0.9,0.9,0.9}{
            \makecell[{{p{\linewidth}}}]{
                \texttt{\tiny{[GM$\langle$P2]}}
                \texttt{Answer: 1} \\
            }
        }
    }
     \\ \\

    \theutterance \stepcounter{utterance}  
    & & & \multicolumn{2}{p{0.3\linewidth}}{
        \cellcolor[rgb]{0.9,0.9,0.9}{
            \makecell[{{p{\linewidth}}}]{
                \texttt{\tiny{[GM$|$GM]}}
                \texttt{1} \\
            }
        }
    }
    & & \\ \\

    \theutterance \stepcounter{utterance}  
    & \multicolumn{4}{p{0.6\linewidth}}{
        \cellcolor[rgb]{0.95,0.95,0.95}{
            \makecell[{{p{\linewidth}}}]{
                \texttt{\tiny{[P2$\rightarrow$P1]}}
                \texttt{Post{-}episode Reflexion (not sent to GM):} \\
\texttt{New plan: In my previous attempt, I focused on the rows of X's, which are present in all three grids. However, I failed to notice that the target grid has a unique feature {-} a row of X's in the bottom row, but also a column of X's in the middle column. I will examine the columns more closely to find a distinguishing characteristic. I will look for a column that has a unique pattern of X's and □'s, or a specific position where the X's are located. For example, I might describe the grid as "the grid that has a column of X's in the middle column and a row of five X's in the bottom row."} \\
            }
        }
    }
    & & \\ \\

    \theutterance \stepcounter{utterance}  
    & & & \multicolumn{4}{p{0.6\linewidth}}{
        \cellcolor[rgb]{0.95,0.95,0.95}{
            \makecell[{{p{\linewidth}}}]{
                \texttt{\tiny{[P1$\rightarrow$P2]}}
                \texttt{Post{-}episode Reflexion (not sent to GM):} \\
\texttt{New plan: I failed to match the given expression because I didn't carefully read it. The expression states that the grid should have a row of five X's in the bottom row. However, I selected the first grid, which only has a row of five X's in the top row. Next time, I will carefully read the expression and look for the additional condition of a row of five X's in the bottom row. I will select the grid that meets both conditions. In this case, I should have noticed that the first grid does not meet the condition of having a row of five X's in the bottom row, and the second and third grids do not meet the condition of having a row of five X's in the top row. Therefore, I should have selected the grid that meets both conditions, which is actually the third grid.} \\
            }
        }
    }
     \\ \\

\end{supertabular}
}

\end{document}
